\section*{Part III: Positioning: Choosing the Meaning of the Comparison}

\subsection*{23. Positioning Is Not Your Tagline}

When someone asks, ``How should we position the business?'', a team may begin designing brand colors or revising the homepage headline. That is usually too late in the reasoning process. A headline expresses a position; it does not discover one. The actual positioning decision determines \textbf{which market comparison the buyer makes and why the offer's real attributes matter inside that comparison.}

Imagine selling software that transcribes meetings. Within the generic meeting-notes category, competitors have adequate transcription and low prices. If the product also enables accountable approvals for regulated client records, the more meaningful frame is \emph{reviewable meeting-record workflow}. That frame changes the comparison set from consumer transcription apps to manual note review, compliance workflows, and the personnel hours required to approve records. The vendor must genuinely provide the controls and integration implied by the new frame. Renaming a generic transcription tool cannot make it a compliant records system.

April Dunford identifies 5 connected components: competitive alternatives, differentiated attributes, customer value, best-fit customers, and market category [R2]. The most useful lesson is that a product may fit more than 1 plausible category, but \textbf{not all frames make the same strengths visible or recruit the same buyer.} Positioning is a set of commercial choices, not a fill-in-the-blanks sentence.

\subsection*{24. Frame Selection Alters Perceived Value Directly}

Consider a 90-minute technical review with a senior engineer. Presented as an \emph{AI workshop}, it competes with inexpensive webinars. Presented as a \emph{production AI design review}, it competes with internal architectural scrutiny, paid consultation, and the cost of releasing an unreliable system. The underlying activity overlaps; the buyer's job, expectations, price reference, and proof obligations differ radically.

This does \emph{not} mean words create an unlimited premium. If the reviewer merely lectures about tools, the design-review label is deceptive. But a truthful frame makes an otherwise overlooked benefit easier to evaluate. Founders sometimes lose because they have the right capability and the wrong comparison. Conversely, founders sometimes persuade customers using a promising frame and then injure trust by delivering an ordinary product.

A good positioning frame must satisfy 4 tests: the customer recognizes the type of problem; the claimed differentiators are factual; the frame invites the right competitive set; and the company can deliver everything the label implies.

\subsection*{25. 3 Bases for a Strategic Position}

Porter's strategy framework distinguishes positioning by product or service variety, by customer needs, and by access method [R1]:

\begin{itemize}
    \item \textbf{Variety-based:} We do 1 type of task exceptionally well. A consultancy may specialize only in model evaluation and deployment readiness rather than generic AI transformation. A repair firm may specialize in 1 narrow machine failure.
    \item \textbf{Needs-based:} We serve an identifiable group comprehensively or unusually well. An accounting platform for independent medical practices may bundle regulatory workflows, revenue reporting, and reconciliations because the users share a distinctive bundle of needs.
    \item \textbf{Access-based:} We reach or serve customers whose location, scale, language, timing, or acquisition requirements differ. A supplier might build an economical local replenishment route where large national distributors are structurally inefficient.
\end{itemize}

These are not mutually exclusive. A useful market position often combines them. The important discipline is that each must imply operational activities, not merely a demographic description.

\subsection*{26. Start with Real Alternatives, Not Product Features}

A useful positioning session begins by suspending the founder's preferred category. Ask a set of recent buyers: What would you have done without us? Which actual products or approaches did you compare? What qualified those options for consideration? Which alternative did you select, and why?

For a new provider of rigorous AI education, the relevant alternatives include inexpensive self-paced content, a university-linked certificate, technical coaching, an open-source project, employer-funded internal training, doing the task at work, and delaying further study. Some are direct substitutes for content; others are substitutes for assessed skill development. They cannot all be beaten by the same promise.

Build an alternative map with buyer job, price or effort, strengths, drawbacks, switching friction, credibility, and approval requirement. Include the status quo explicitly. Differentiation that ignores inaction is particularly fragile in markets with optional purchases and high time commitments.

\subsection*{27. The Alternative Determines the Evidence Burden}

Suppose an instructor tells an experienced engineer that the course offers ``high-quality videos.'' The engineer can obtain videos for free. That is a weak claim relative to the alternative. Now suppose the instructor offers timed review of the learner's actual architecture and a demanding technical defense. The buyer finds that harder to reproduce through solo study.

Compared with a competent senior colleague willing to mentor the engineer, the superiority may disappear. This is why the honest competitor map must include strong, inconvenient rivals. A founder who can articulate the best alternative is more likely to discover which differences are genuinely scarce.

Different alternatives require different proof. If competing with a textbook, demonstrate improved practice and feedback. If competing with a named prestigious program, demonstrate delivery outcomes while conceding weaker credential recognition. If competing with internal consultancy, show faster or higher-quality diagnosis net of organizational costs.

\subsection*{28. Translate Attributes into Value and Consequences}

A feature is what the offering contains. A capability is what the organization can do. A benefit is how that improves the buyer's situation. A consequence is why the benefit matters at this time:

\begin{itemize}
    \item \textbf{Feature:} 6 architecture review sessions.
    \item \textbf{Capability:} Reviewers identify significant design and evaluation failures in authentic projects.
    \item \textbf{Immediate benefit:} Learners receive timely correction on consequential decisions.
    \item \textbf{Consequence:} They independently defend stronger technical choices, avoiding expensive rework in production.
    \item \textbf{Proof:} Published review standards, sample critiques, independent project assessment, and observed learner performance.
\end{itemize}

The weakest marketing jumps from ``6 sessions'' directly to ``change your career.'' The strongest stops at the farthest implication the evidence can support. In B2B sales, James Anderson, James Narus, and Wouter van Rossum emphasize the necessity of substantiating customer value rather than relying on unsupported benefit assertions [R5].
